% !TeX program = pdflatex
\documentclass[11pt,a4paper]{amsart}

\newcommand{\notelanguage}{finnish}
% Shared preamble for course notes and exercise sets.

\usepackage[T1]{fontenc}
\usepackage{lmodern}
\usepackage[english]{babel}

\usepackage[a4paper,margin=30mm]{geometry}
\usepackage{microtype}
\usepackage{graphicx}
\usepackage{enumitem}

\usepackage{amsmath}
\usepackage{amssymb}
\usepackage{amsthm}
\usepackage{mathtools}
\usepackage{aliascnt}

\usepackage[hidelinks,pdfusetitle]{hyperref}

\numberwithin{equation}{section}
\setlist{itemsep=0.25em,topsep=0.5em}

\theoremstyle{plain}
\newtheorem{theorem}{Theorem}[section]
\newaliascnt{lemma}{theorem}
\newtheorem{lemma}[lemma]{Lemma}
\aliascntresetthe{lemma}
\newaliascnt{proposition}{theorem}
\newtheorem{proposition}[proposition]{Proposition}
\aliascntresetthe{proposition}
\newaliascnt{corollary}{theorem}
\newtheorem{corollary}[corollary]{Corollary}
\aliascntresetthe{corollary}

\theoremstyle{definition}
\newaliascnt{definition}{theorem}
\newtheorem{definition}[definition]{Definition}
\aliascntresetthe{definition}
\newaliascnt{example}{theorem}
\newtheorem{example}[example]{Example}
\aliascntresetthe{example}
\newaliascnt{problem}{theorem}
\newtheorem{problem}[problem]{Problem}
\aliascntresetthe{problem}

\theoremstyle{remark}
\newaliascnt{remark}{theorem}
\newtheorem{remark}[remark]{Remark}
\aliascntresetthe{remark}

\usepackage[nameinlink,noabbrev]{cleveref}

\crefname{theorem}{theorem}{theorems}
\crefname{lemma}{lemma}{lemmas}
\crefname{proposition}{proposition}{propositions}
\crefname{corollary}{corollary}{corollaries}
\crefname{definition}{definition}{definitions}
\crefname{example}{example}{examples}
\crefname{problem}{problem}{problems}
\crefname{remark}{remark}{remarks}

\DeclareMathOperator{\im}{im}
\DeclareMathOperator{\rank}{rank}
\DeclarePairedDelimiter{\abs}{\lvert}{\rvert}
\DeclarePairedDelimiter{\norm}{\lVert}{\rVert}
\DeclarePairedDelimiter{\set}{\lbrace}{\rbrace}


\title{Johdatus logiikkaan I\\Harjoituksen 4 ratkaisut}
\author{}
\date{}

\begin{document}

\exerciseheading{Johdatus logiikkaan I}{Harjoituksen 4 ratkaisut}

\begin{exercise}
  Muunna seuraavat lauseet klausuulimuotoon:
\begin{enumerate}[label=\alph*), nosep]
  \item \((p_0 \lor p_1) \land (p_2 \rightarrow \neg p_0)\)
  \item \(\neg(p_0 \rightarrow p_2) \land \neg(p_1 \leftrightarrow p_2)\)
  \item \((p_0 \lor p_1) \rightarrow (p_2 \lor p_3)\)
  \item \((p_0 \land p_1) \rightarrow (p_2 \land p_3)\)
  \item \((\neg p_1 \land p_0) \lor (\neg p_2 \land p_0)\)
\end{enumerate}
\end{exercise}

\begin{solution}
Muunnetaan lauseet ensin ekvivalenttiin konjunktiiviseen normaalimuotoon.
Jokainen disjunktio kirjoitetaan literaalien joukkona ja konjunktio
näiden klausuulien joukkona.
\begin{enumerate}[label=\alph*)]
\item Poistetaan implikaatio:
  \[
    (p_0 \lor p_1) \land (p_2 \to \neg p_0)
    \iff (p_0 \lor p_1) \land (\neg p_2 \lor \neg p_0).
  \]
  Klausuulimuoto on
  \[
    \big\{\{p_0,p_1\},\{\neg p_2,\neg p_0\}\big\}.
  \]

\item Implikaation negaatio saadaan De Morganin kaavalla:
  \[
    \neg(p_0 \to p_2)
    \iff \neg(\neg p_0 \lor p_2)
    \iff p_0 \land \neg p_2.
  \]
  Ekvivalenssin negaatiolle saadaan vastaavasti
  \[
    \begin{aligned}
      \neg(p_1 \leftrightarrow p_2)
      &\iff \neg\big[(\neg p_1 \lor p_2)\land(\neg p_2 \lor p_1)\big] \\
      &\iff \neg(\neg p_1 \lor p_2)\lor\neg(\neg p_2 \lor p_1) \\
      &\iff (p_1\land\neg p_2)\lor(p_2\land\neg p_1) \\
      &\iff (p_1\lor p_2)\land(p_1\lor\neg p_1) \\
      &\qquad\land(\neg p_2\lor p_2)\land(\neg p_2\lor\neg p_1) \\
      &\iff (p_1\lor p_2)\land(\neg p_1\lor\neg p_2).
    \end{aligned}
  \]
  Viimeisessä vaiheessa poistettiin kaksi tautologista konjunktia.
  Yhdistämällä saadaan CNF
  \[
    p_0\land\neg p_2\land(p_1\lor p_2)\land(\neg p_1\lor\neg p_2),
  \]
  joten klausuulimuoto on
  \[
    \big\{\{p_0\},\{\neg p_2\},\{p_1,p_2\},\{\neg p_1,\neg p_2\}\big\}.
  \]

\item Poistetaan implikaatio ja sovelletaan De Morganin kaavaa sekä
  distributiivisuutta:
  \[
    \begin{aligned}
      (p_0\lor p_1)\to(p_2\lor p_3)
      &\iff \neg(p_0\lor p_1)\lor(p_2\lor p_3) \\
      &\iff (\neg p_0\land\neg p_1)\lor(p_2\lor p_3) \\
      &\iff (\neg p_0\lor p_2\lor p_3)
      \land(\neg p_1\lor p_2\lor p_3).
    \end{aligned}
  \]
  Klausuulimuoto on
  \[
    \big\{\{\neg p_0,p_2,p_3\},\{\neg p_1,p_2,p_3\}\big\}.
  \]

\item Vastaavasti
  \[
    \begin{aligned}
      (p_0\land p_1)\to(p_2\land p_3)
      &\iff \neg(p_0\land p_1)\lor(p_2\land p_3) \\
      &\iff (\neg p_0\lor\neg p_1)\lor(p_2\land p_3) \\
      &\iff (\neg p_0\lor\neg p_1\lor p_2)
      \land(\neg p_0\lor\neg p_1\lor p_3).
    \end{aligned}
  \]
  Klausuulimuoto on
  \[
    \big\{\{\neg p_0,\neg p_1,p_2\},\{\neg p_0,\neg p_1,p_3\}\big\}.
  \]

\item Otetaan yhteinen konjunkti \(p_0\) distributiivisuudella ulos:
  \[
    (\neg p_1\land p_0)\lor(\neg p_2\land p_0)
    \iff (\neg p_1\lor\neg p_2)\land p_0.
  \]
  Klausuulimuoto on
  \[
    \big\{\{\neg p_1,\neg p_2\},\{p_0\}\big\}.
  \]
\end{enumerate}
\end{solution}

\begin{exercise}
Minkä klausuulin resoluutiosääntö tuottaa seuraavista klausuulipareista?
\begin{enumerate}[label=\alph*), nosep]
\item \(\{\neg p_0, p_1, p_2\}\) ja \(\{p_0, p_1, p_3\}\)
\item \(\{p_1, \neg p_1, \neg p_2\}\) ja \(\{\neg p_1, p_3, \neg p_4\}\)
\item \(\{p_0, \neg p_1, \neg p_2\}\) ja \(\{\neg p_0, p_1, p_3\}\)
\item \(\{\neg p_0, p_1, p_2\}\) ja \(\{p_2, \neg p_3, p_4\}\)
\end{enumerate}
Onko vastaus aina yksikäsitteinen?
\end{exercise}

\begin{solution}
Valitaan klausuuleista vastakkaiset literaalit \(p_i\) ja \(\neg p_i\),
poistetaan valittu literaali kummastakin klausuulista ja otetaan jäljelle
jäävien joukkojen yhdiste. Kerralla poistetaan vain yksi vastakkainen pari.
\begin{enumerate}[label=\alph*)]
\item Resolvoidaan propositiosymbolin \(p_0\) suhteen:
\[
\{p_1,p_2\}\cup\{p_1,p_3\}=\{p_1,p_2,p_3\}.
\]
Joukossa sama literaali esiintyy vain kerran.

\item Poistetaan \(p_1\) ensimmäisestä klausuulista ja \(\neg p_1\)
toisesta klausuulista:
\[
\{\neg p_1,\neg p_2\}\cup\{p_3,\neg p_4\}
=\{\neg p_1,\neg p_2,p_3,\neg p_4\}.
\]
Ensimmäisen klausuulin literaali \(\neg p_1\) jää mukaan.

\item Voidaan resolvoida joko \(p_0\):n tai \(p_1\):n suhteen.
Propositiosymbolin \(p_0\) suhteen saadaan
\[
\{\neg p_1,\neg p_2\}\cup\{p_1,p_3\}
=\{\neg p_1,\neg p_2,p_1,p_3\}.
\]
Propositiosymbolin \(p_1\) suhteen saadaan
\[
\{p_0,\neg p_2\}\cup\{\neg p_0,p_3\}
=\{p_0,\neg p_2,\neg p_0,p_3\}.
\]
Molemmat tulokset ovat tautologisia klausuuleja, mutta kelvollisia
resoluutiosäännön tuloksia. Molempia vastakkaisia pareja ei saa poistaa
samalla askeleella.

\item Resoluutiosääntöä ei voi soveltaa näihin kahteen klausuuliin,
sillä niissä ei ole vastakkaista literaaliparia. Yhteinen literaali
\(p_2\) ei riitä: toisessa klausuulissa pitäisi esiintyä \(\neg p_2\).
\end{enumerate}
Vastaus ei siis aina ole yksikäsitteinen: kohdassa (c) eri
propositiosymbolin valinta antaa eri klausuulin.
\end{solution}

\begin{exercise}
Päättele resoluutiolla \(\emptyset\) seuraavista klausuulijoukoista:
\begin{enumerate}[label=\alph*), nosep]
\item \(\{\{\neg p_0\}, \{\neg p_1\}, \{p_0, p_2\}, \{p_1, \neg p_2\}\}\)
\item \(\{\{\neg p_0, \neg p_1\}, \{p_0, p_2\}, \{p_0, \neg p_2\},
\{\neg p_0, p_1\}\}\)
\item \(\{\{p_1, p_2, \neg p_3\}, \{\neg p_2\}, \{\neg p_1, p_2\},
\{p_1, p_3\}\}\)
\end{enumerate}
\end{exercise}

\begin{solution}
Osoitetaan klausuulijoukot ristiriitaisiksi resoluutiolla:
\begin{enumerate}[label=\alph*)]
\item
\begin{enumerate}[label=\arabic*., nosep]
\item \(\{\neg p_0\}\)
\item \(\{\neg p_1\}\)
\item \(\{p_0,p_2\}\)
\item \(\{p_1,\neg p_2\}\)
\item \(\{p_2\}\) (rivit 1 ja 3)
\item \(\{\neg p_2\}\) (rivit 2 ja 4)
\item \(\emptyset\) (rivit 5 ja 6)
\end{enumerate}

\item
\begin{enumerate}[label=\arabic*., nosep]
\item \(\{\neg p_0,\neg p_1\}\)
\item \(\{p_0,p_2\}\)
\item \(\{p_0,\neg p_2\}\)
\item \(\{\neg p_0,p_1\}\)
\item \(\{p_0\}\) (rivit 2 ja 3)
\item \(\{\neg p_0\}\) (rivit 1 ja 4)
\item \(\emptyset\) (rivit 5 ja 6)
\end{enumerate}

\item
\begin{enumerate}[label=\arabic*., nosep]
\item \(\{p_1,p_2,\neg p_3\}\)
\item \(\{\neg p_2\}\)
\item \(\{\neg p_1,p_2\}\)
\item \(\{p_1,p_3\}\)
\item \(\{p_1,\neg p_3\}\) (rivit 1 ja 2)
\item \(\{\neg p_1\}\) (rivit 2 ja 3)
\item \(\{p_1\}\) (rivit 4 ja 5)
\item \(\emptyset\) (rivit 6 ja 7)
\end{enumerate}
\end{enumerate}
\end{solution}

\begin{exercise}
Osoita resoluution avulla, että \((p_0 \lor p_2) \rightarrow (p_1 \lor p_3)\)
on oletusten \((p_0 \rightarrow p_1)\) ja \((p_2 \rightarrow p_3)\)
looginen seuraus.
\end{exercise}

\begin{solution}
Muunnetaan oletukset klausuulimuotoon:
\[
p_0\to p_1\iff\neg p_0\lor p_1,
\qquad
p_2\to p_3\iff\neg p_2\lor p_3.
\]
Lisätään johtopäätöksen negaatio, jolle saadaan
\[
\begin{aligned}
\neg\big[(p_0\lor p_2)\to(p_1\lor p_3)\big]
&\iff (p_0\lor p_2)\land\neg(p_1\lor p_3)\\
&\iff (p_0\lor p_2)\land\neg p_1\land\neg p_3.
\end{aligned}
\]
Osoitetaan näin saatu klausuulijoukko ristiriitaiseksi resoluutiolla:
\begin{enumerate}[label=\arabic*., nosep]
\item \(\{\neg p_0,p_1\}\)
\item \(\{\neg p_2,p_3\}\)
\item \(\{p_0,p_2\}\)
\item \(\{\neg p_1\}\)
\item \(\{\neg p_3\}\)
\item \(\{\neg p_0\}\) (rivit 1 ja 4)
\item \(\{p_2\}\) (rivit 3 ja 6)
\item \(\{p_3\}\) (rivit 2 ja 7)
\item \(\emptyset\) (rivit 5 ja 8)
\end{enumerate}
Oletukset yhdessä johtopäätöksen negaation kanssa ovat siis ristiriitaiset.
Näin ollen \((p_0\lor p_2)\to(p_1\lor p_3)\) on annettujen oletusten
looginen seuraus.
\end{solution}

\begin{exercise}
Osoita resoluution avulla, että joukko
\[
\{\neg p_0 \rightarrow p_1, p_1 \rightarrow p_0,
p_0 \rightarrow (p_2 \land p_3), \neg p_0 \lor \neg p_2 \lor \neg p_3\}
\]
ei ole toteutuva. Perustuuko resoluution toimivuus todistusmenetelmänä tässä
sen eheyteen, kumoustäydellisyyteen tai molempiin? (Eli mitä eheys ja
kumoustäydellisyys tässä antavat?)
\end{exercise}

\begin{solution}
Muunnetaan implikaatiot konjunktiiviseen normaalimuotoon:
\[
\begin{aligned}
\neg p_0\to p_1 &\iff p_0\lor p_1,\\
p_1\to p_0 &\iff \neg p_1\lor p_0,\\
p_0\to(p_2\land p_3)
&\iff \neg p_0\lor(p_2\land p_3)\\
&\iff (\neg p_0\lor p_2)\land(\neg p_0\lor p_3).
\end{aligned}
\]
Viimeinen oletus on jo literaalien disjunktio. Osoitetaan saatu
klausuulijoukko ristiriitaiseksi resoluutiolla:
\begin{enumerate}[label=\arabic*., nosep]
\item \(\{p_0,p_1\}\)
\item \(\{\neg p_1,p_0\}\)
\item \(\{\neg p_0,p_2\}\)
\item \(\{\neg p_0,p_3\}\)
\item \(\{\neg p_0,\neg p_2,\neg p_3\}\)
\item \(\{p_0\}\) (rivit 1 ja 2)
\item \(\{p_2\}\) (rivit 3 ja 6)
\item \(\{p_3\}\) (rivit 4 ja 6)
\item \(\{\neg p_2,\neg p_3\}\) (rivit 5 ja 6)
\item \(\{\neg p_3\}\) (rivit 7 ja 9)
\item \(\emptyset\) (rivit 8 ja 10)
\end{enumerate}
Resoluution \emph{eheyden} nojalla johdettu tyhjä klausuuli on oletusten
looginen seuraus. Koska tyhjä klausuuli on aina epätosi, mikään
totuusjakauma ei voi toteuttaa kaikkia oletuksia. Alkuperäinen joukko
ei siis ole toteutuva.

\emph{Kumoustäydellisyys} puolestaan takaa, että jokaisesta
ristiriitaisesta klausuulijoukosta voidaan johtaa tyhjä klausuuli
resoluutiolla. Se takaa siis tällaisen todistuksen olemassaolon, kun
joukko on ristiriitainen. Yllä löydetystä todistuksesta tehtävä päätelmä
joukon toteutumattomuudesta perustuu eheyteen; kumoustäydellisyyttä ei
tarvita tämän yksittäisen todistuksen johtopäätöksen perustelemiseen.
\end{solution}

\begin{exercise}
Tarkastellaan seuraavaa ”resoluutiotodistusta”, joka on osoittavinaan, että
\(p_2\) olisi oletusten \(p_0 \rightarrow (p_1 \lor p_2)\) ja
\(p_1 \rightarrow p_0\) looginen seuraus.
\begin{enumerate}[label=\arabic*., nosep]
\item \(\{\neg p_0, p_1, p_2\}\)
\item \(\{\neg p_1, p_0\}\)
\item \(\{p_2\}\)
\item \(\{\neg p_2\}\)
\item \(\emptyset\)
\end{enumerate}
Miten tiedämme jo ennen todistukseen perehtymistä, että siinä on oltava
virhe? Mitä ”resoluutiotodistuksessa” on tehty, ja missä on tapahtunut virhe?
\end{exercise}

\begin{solution}
Valitaan totuusjakauma \(v\), jolla
\[
v(p_0)=1,\qquad v(p_1)=1,\qquad v(p_2)=0.
\]
Tällöin molemmat oletukset \(p_0\to(p_1\lor p_2)\) ja
\(p_1\to p_0\) ovat tosia, mutta väitetty johtopäätös \(p_2\) on
epätosi. Siis \(p_2\) ei ole oletusten looginen seuraus. Resoluution
eheyden perusteella esitetyssä todistuksessa on oltava virhe.

Riveillä 1 ja 2 oletukset on muunnettu oikein klausuulimuotoon.
Virhe tapahtuu rivillä 3: molemmat vastakkaiset literaaliparit
\(p_0,\neg p_0\) ja \(p_1,\neg p_1\) on poistettu samalla askeleella.
Resoluutiosäännöllä poistetaan vain yksi vastakkainen pari kerrallaan.
Resoluutiolla \(p_0\):n suhteen saataisiin
\[
\{p_1,p_2\}\cup\{\neg p_1\}=\{p_1,\neg p_1,p_2\},
\]
ja resoluutiolla \(p_1\):n suhteen saataisiin
\[
\{\neg p_0,p_2\}\cup\{p_0\}=\{\neg p_0,p_0,p_2\}.
\]
Kummassakaan tapauksessa tulos ei ole \(\{p_2\}\). Saadut klausuulit
ovat tautologisia, eikä niiden sisällä olevaa vastakkaista literaaliparia
voi poistaa.

Rivillä 4 on lisätty väitetyn johtopäätöksen negaatio, kuten
resoluutiomenetelmässä kuuluukin. Rivin 5 resoluutio riveistä 3 ja 4
olisi kelvollinen, jos rivi 3 olisi johdettu oikein. Todistuksen virhe
on siis nimenomaan rivin 3 johtamisessa.
\end{solution}

\begin{exercise}
Lemman 6.7 todistuksessa klausuulijoukko \(C\) jaetaan neljään osaan.
Miten käy, jos jokin osa on tyhjä? Osoita, että kumoustäydellisyys pätee
näissäkin tapauksissa.
\end{exercise}

\begin{solution}
Olkoon \(C\) ristiriitainen ja käytetään lemman merkintöjä
\[
C'=C_{00}\cup(C_{01}\times_R C_{10}).
\]
Jos \(C_{11}=\emptyset\), todistus ei muutu, sillä tämän osan klausuulit
ovat tautologisia. Jos \(C_{00}=\emptyset\), lemman todistus
\(C'\):n ristiriitaisuudesta toimii sellaisenaan: ehto, että totuusjakauma
toteuttaa kaikki \(C_{00}\):n klausuulit, pätee tyhjälle joukolle aina.

Jos \(C_{01}=\emptyset\), niin \(C'=C_{00}\). Jos \(C_{00}\) olisi
toteutuva, sitä toteuttava totuusjakauma voitaisiin laajentaa asettamalla
\(p_i=1\). Tällöin myös kaikki \(C_{10}\):n ja \(C_{11}\):n klausuulit
olisivat tosia, joten \(C\) olisi toteutuva, vastoin oletusta.
Siis \(C'\) on ristiriitainen. Jos \(C_{10}=\emptyset\), sama perustelu
toimii asettamalla \(p_i=0\).

Nämä perustelut kattavat myös usean osan samanaikaisen tyhjyyden.
Joukossa \(C'\) esiintyy vähemmän propositiosymboleja kuin joukossa
\(C\), joten induktio-oletuksella \(C'\vdash_R\emptyset\).
Koska \(C'\):n klausuulit ovat joko \(C\):n klausuuleja tai niistä
resoluutiolla johdettuja, myös \(C\vdash_R\emptyset\).
\end{solution}

\begin{exercise}
Propositiologiikan resoluutio pysähtyy aina, eli jossain vaiheessa ei saada
pääteltyä uusia klausuuleja. Etsi yläraja resoluution pituudelle, jos
oletukset koostuvat \(m\) klausuulista, joissa kussakin esiintyy korkeintaan
\(n\) literaalia. Huomaa: mikä tahansa yläraja (perusteluineen) kelpaa,
ei tarvitse etsiä optimaalista ylärajaa.
\end{exercise}

\begin{solution}
Olkoon \(L\) kaikkien oletusklausuuleissa esiintyvien literaalien joukko.
Koska oletuksia on \(m\) ja kussakin on korkeintaan \(n\) literaalia,
niin \(|L|\le mn\). Resoluutio ei tuota uusia literaaleja: resolventti
muodostetaan poistamalla vastakkaiset literaalit ja ottamalla jäljelle
jäävien joukkojen yhdiste. Siksi jokainen johdettu klausuuli on
joukon \(L\) osajoukko.

Mahdollisia klausuuleja on siis korkeintaan
\[
|\mathcal{P}(L)|=2^{|L|}\le 2^{mn},
\]
missä myös tyhjä klausuuli on mukana. Kun jokaisella askeleella lisätään
aidosti uusi klausuuli, askelia voi näin ollen olla korkeintaan
\(2^{mn}\). Jos myös alkuperäiset \(m\) oletusriviä lasketaan mukaan,
resoluutiotodistuksen pituudelle kelpaa yläraja \(m+2^{mn}\).

Tässä samoja klausuuleja ei johdeta uudelleen: jos toistoja sallittaisiin,
todistusta voisi pitkittää mielivaltaisesti. Pysähtyminen tarkoittaa,
että lopulta uusia klausuuleja ei enää saada.
\end{solution}

\begin{exercise}
Lemma 6.7 antaa suoraviivaisen menetelmän tuottaa \(\emptyset\)
ristiriitaisesta klausuulijoukosta. Kuinka pitkä resoluutiotodistuksesta
tulee, jos sovellat lemman menetelmää klausuulijoukkoon
\[
\{\{\neg p_0, p_2\}, \{\neg p_1, p_2\}, \{\neg p_2, p_3\},
\{p_0, p_1, \neg p_3\}, \{p_1, p_2, p_3\}, \{\neg p_2, \neg p_3\}\}
\]
Huomaa, että valittu propositiosymboleiden käsittelyjärjestys vaikuttaa
vastaukseen. Kuinka lyhyen resoluutiotodistuksen keksit, jos et noudata
lemman menetelmää suoraan, vaan optimoit valintoja?
\end{exercise}

\begin{solution}
Numeroidaan alkuperäiset klausuulit annetussa järjestyksessä riveiksi
1--6. Käytetään lemman menetelmässä käsittelyjärjestystä
\(p_0,p_1,p_2,p_3\). Kussakin vaiheessa muodostetaan kaikki valitun
symbolin suhteen saadut resolventit ja säilytetään klausuulit, joissa
symbolia ei esiinny kummallakaan merkillä.

Symbolin \(p_0\) eliminointi tuottaa riveistä 1 ja 4 klausuulin
\(\{p_1,p_2,\neg p_3\}\). Uusi klausuulijoukko on siis
\[
C^{(1)}=\bigl\{
\{\neg p_1,p_2\},\{\neg p_2,p_3\},\{p_1,p_2,p_3\},
\{\neg p_2,\neg p_3\},\{p_1,p_2,\neg p_3\}
\bigr\}.
\]
Seuraavaksi \(\{\neg p_1,p_2\}\) resolvoidaan kummankin
\(p_1\):n sisältävän klausuulin kanssa. Saadaan
\[
C^{(2)}=\bigl\{
\{\neg p_2,p_3\},\{\neg p_2,\neg p_3\},
\{p_2,p_3\},\{p_2,\neg p_3\}
\bigr\}.
\]
Symbolin \(p_2\) suhteen on neljä klausuuliparia. Niistä saadaan
kolme eri resolventtia, sillä kaksi paria antaa saman tautologisen
klausuulin:
\[
C^{(3)}=\bigl\{\{p_3\},\{p_3,\neg p_3\},\{\neg p_3\}\bigr\}.
\]
Viimeisessä vaiheessa tautologinen klausuuli kuuluu osaan \(C_{11}\)
ja jätetään pois. Klausuuleista \(\{p_3\}\) ja \(\{\neg p_3\}\)
saadaan \(\emptyset\).

Koko resoluutiotodistus on seuraava; sama klausuuli kirjoitetaan vain
kerran:
\begin{enumerate}[label=\arabic*., nosep]
\item \(\{\neg p_0,p_2\}\)
\item \(\{\neg p_1,p_2\}\)
\item \(\{\neg p_2,p_3\}\)
\item \(\{p_0,p_1,\neg p_3\}\)
\item \(\{p_1,p_2,p_3\}\)
\item \(\{\neg p_2,\neg p_3\}\)
\item \(\{p_1,p_2,\neg p_3\}\) (rivit 1 ja 4, symboli \(p_0\))
\item \(\{p_2,p_3\}\) (rivit 2 ja 5, symboli \(p_1\))
\item \(\{p_2,\neg p_3\}\) (rivit 2 ja 7, symboli \(p_1\))
\item \(\{p_3\}\) (rivit 3 ja 8, symboli \(p_2\))
\item \(\{p_3,\neg p_3\}\) (rivit 6 ja 8, myös rivit 3 ja 9,
symboli \(p_2\))
\item \(\{\neg p_3\}\) (rivit 6 ja 9, symboli \(p_2\))
\item \(\emptyset\) (rivit 10 ja 12, symboli \(p_3\))
\end{enumerate}
Näin saadaan \(7\) resoluutioaskelta eli \(13\) riviä oletukset mukaan
lukien. Jos myös saman tautologisen klausuulin toinen johtaminen
kirjoitetaan erikseen, saadaan \(8\) askelta eli \(14\) riviä.

Lyhyempi todistus saadaan käyttämällä samoja oletusrivejä 1--6 ja
valitsemalla jatko seuraavasti:
\begin{enumerate}[label=\arabic*., start=7, nosep]
\item \(\{p_1,p_2,\neg p_3\}\) (rivit 1 ja 4, symboli \(p_0\))
\item \(\{p_1,p_2\}\) (rivit 5 ja 7, symboli \(p_3\))
\item \(\{p_2\}\) (rivit 2 ja 8, symboli \(p_1\))
\item \(\{\neg p_2\}\) (rivit 3 ja 6, symboli \(p_3\))
\item \(\emptyset\) (rivit 9 ja 10, symboli \(p_2\))
\end{enumerate}
Tässä on vain \(5\) resoluutioaskelta eli \(11\) riviä oletukset mukaan
lukien.
\end{solution}

\begin{exercise}
Resoluutiosäännön voi määritellä kahdessa muodossa:
\[
\frac{C_1 \cup \{p_i\} \text{ ja } C_2 \cup \{\neg p_i\}}{C_1 \cup C_2}
\qquad\text{tai}\qquad
\frac{C_1 \mathbin{\dot\cup} \{p_i\} \text{ ja }
C_2 \mathbin{\dot\cup} \{\neg p_i\}}{C_1 \cup C_2}
\]
Vasemmanpuoleisessa määritelmässä sallitaan mielivaltaiset yhdisteet,
eli myös esim. sellaiset, jossa \(p_i \in C_1\). Oikeanpuoleisessa
tarkastellaan \emph{erillisiä yhdisteitä}, eli esim.
\(C_1 \cap \{p_i\} = \emptyset\). Miten määritelmän muuttaminen
vaikuttaa resoluution käyttöön?
\end{exercise}

\begin{solution}
Erillisten yhdisteiden määritelmässä valittu literaali on poistettava
kummastakin lähtöklausuulista. Tavallisten yhdisteiden määritelmässä
sen voi myös jättää joukkoon \(C_1\) tai \(C_2\), jolloin se jää mukaan
lopputulokseen.

Esimerkiksi klausuuleista \(\{p_0\}\) ja \(\{\neg p_0,p_1\}\)
saadaan erillisten yhdisteiden säännöllä vain \(\{p_1\}\).
Tavallisten yhdisteiden säännöllä mahdolliset tulokset ovat
\[
\{p_1\},\quad \{p_0,p_1\},\quad \{\neg p_0,p_1\},\quad
\{p_0,\neg p_0,p_1\}.
\]
Esimerkiksi tulos \(\{p_0,p_1\}\) saadaan valitsemalla
\(C_1=\{p_0\}\) ja \(C_2=\{p_1\}\), sillä
\(\{p_0\}\cup\{p_0\}=\{p_0\}\).

Tavallisten yhdisteiden sääntö antaa siis enemmän valinnanvaraa,
mutta ylimääräiset tulokset ovat tavallisen resolventin heikennyksiä.
Käytännössä resoluutiota voi edelleen käyttää kuten ennenkin:
poistetaan valittu vastakkainen literaalipari ja yhdistetään jäljelle
jäävät literaalit. Erillisten yhdisteiden määritelmä tekee tämän
poistamisen pakolliseksi.
\end{solution}

\end{document}
